\documentclass[10pt,a4paper]{amsart}
\usepackage[margin=19mm]{geometry}
\usepackage{amsmath,amssymb,mathtools}
\usepackage[scr=rsfs,cal=euler]{mathalfa}
\usepackage{xfrac}
\usepackage[unicode,hidelinks,bookmarks=false]{hyperref}
\newcommand{\ie}{i.e.\ }
\newcommand{\cf}{cf.\ }
\newcommand{\resp}{respectively\ }
\newcommand{\Subsection}{\subsection}
\providecommand{\nobreakdash}{\nobreak\hbox{-}}
\setlength{\emergencystretch}{2em}
\multlinegap0pt
\usepackage{mathtools}
\usepackage{amssymb}
\usepackage{xcolor}
\usepackage{dsfont}
\usepackage{tikz}
\usepackage{bm}
\usepackage[shortlabels]{enumitem}
\newtheorem{corollary}{Corollary}
\newtheorem{theorem}{Theorem}
\def\BB{\mathbb{B}}
\def\EE{\mathbb{E}}
\def\HH{\mathbb{H}}
\def\RR{\mathbb{R}}
\def\SS{\mathbb{S}}
\def\bp{\mathbf{p}}
\newcommand{\bu}{\boldsymbol{u}}
\newcommand{\bw}{\boldsymbol{w}}
\newcommand{\bz}{\boldsymbol{z}}
\def\cD{\mathcal{D}}
\def\cH{\mathcal{H}}
\newcommand{\dint}{\textup{d}}
\newcommand{\inter}{\operatorname{int}}
\newcommand{\vol}{\mathop{\mathrm{vol}}\nolimits}
\title{\bfseries Random hyperbolic polyhedra in horoballs}
\author{Florian Besau, Anna Gusakova, Christoph Th\"ale}
\date{Source: arXiv:2609.10007v1 (CC BY 4.0)}
\begin{document}
\maketitle

\noindent\emph{Source and licence.} \bfseries Random hyperbolic polyhedra in horoballs. Version arXiv:2609.10007v1 (CC BY 4.0), distributed by arXiv under the Creative Commons Attribution 4.0 International licence, http://creativecommons.org/licenses/by/4.0/, as recorded in the arXiv licence metadata for this exact version and shown on the abstract page https://arxiv.org/abs/2609.10007v1. Full licence notice: \texttt{licenses/arxiv-2609.10007-CC-BY-4.0.txt}.\par\smallskip

\noindent\textit{Editorial scope.} Counted unit: the article's theorem thm:AsymptoticsGeneralD (source0.tex lines 1226--1236). The article's complete original proof of it is delivered with it (source0.tex lines 3221--3528, one \texttt{proof} environment of 308 lines, which the article places away from the statement). No mathematics is written, repaired or re-derived: the statement and the proof are the article's own lines, in the article's own notation, and no retained line is substituted or re-flowed.  Retained uncounted as the source-local context the proof needs: lines 577--580; lines 697--700; lines 1215--1222; lines 1812--1834; lines 1841--1848; lines 1861--1872; lines 1885--1887; lines 1904--1928; lines 2827--2830; lines 2832--2835.  Nothing was omitted from any retained span, so the package carries no omission record.  No retained line contains a citation, so the wrapper carries no bibliography and no bibliography entry.  No retained line contains a figure environment, an image-inclusion command or drawing code, so no image is delivered and the package declares no asset dependency.  Editorial formatting only, in addition: an AMS \texttt{amsart} preamble that restates the article's own declarations and macros used by the retained lines, namely the environments \texttt{corollary}, \texttt{theorem} on separate counters, together with the standard packages those lines need.  The retained environments print in this document as Corollary 1, Theorem 1 in order of appearance, whereas the article prints the same items with the numbering of its own full document.  The complete native source of the article as distributed in the arXiv e-print for this exact version is bundled unchanged as \texttt{sources/209866/source0.tex}.
\begin{equation}
	\label{eqn:VolumeCylinder}
	\cH_{\HH^d}^{d-1}(W_\lambda) = (d-1)\, \vol_{\HH^d}\bigl(C_\lambda(W)\bigr)  = \vol_{\RR^{d-1}}\bigl(\tfrac{1}{\lambda} W\bigr).
\end{equation}

\begin{equation}\label{eq:HyperbolicVolume}
	\vol_{\HH^d}(A)
	= \int_A y^{-d}\,\dint v_1\cdots\dint v_{d-1}\,\dint y.
\end{equation}

\begin{corollary}\label{cor:lambda_asymptotics_d2}
	Let $d=2$ and set $\gamma=\lambda$. Then, as $\lambda\downarrow 0$,
	\[
		\xi_1(\lambda,\lambda)
		=1-\pi\lambda+o(\lambda).
	\]
	In particular, $\lim_{\lambda\downarrow 0}\xi_1(\lambda,\lambda)=1$.
\end{corollary}

\begin{theorem}\label{thm:LocalVolume}
	Let $W\subset\RR^{d-1}$ be a bounded Borel set with
	$0<\vol_{\RR^{d-1}}(W)<\infty$. Then, for
	$\gamma,\lambda>0$ and $\ell\in\{0,1\}$,
	\[
		\EE F_{\gamma}^{(\ell)}:=\EE F_{\gamma,\lambda}^{(\ell)}(W)
		=
		\frac{\gamma^d}{2d}
		\int_0^1
		(1-s)^{\frac{d-3}{2}}
		\exp\!\Big(
		-\gamma\frac{\kappa_{d-1}}{d+1}
		s^{\frac{d+1}{2}}
		{}_2F_1\!\Big(
		\tfrac{d+1}{2},\tfrac{d+1}{2};
		\tfrac{d+3}{2};s
		\Big)
		\Big)
		J_{d,\ell}(s)\,\dint s.
	\]
	In particular, both expectations $\EE V_{\gamma}=\EE F_{\gamma}^{(0)}=\EE V_{\gamma,\lambda}(W)$ and $\EE S_{\gamma}=\EE F_{\gamma}^{(1)}=\EE S_{\gamma,\lambda}(W)$ are independent of $W$ and
	$\lambda$.
\end{theorem}

\begin{corollary}\label{cor:Efron}
	For the vertex intensity we have
	\[
		\lambda^{d-1}\, \xi_0(\gamma,\lambda)
		=\frac{\gamma}{(d-1)}\,
		\bigl(1- \EE V_{\gamma}\bigr).
	\]
\end{corollary}

\begin{equation}\label{eq:Gdl}
	G_{d,\ell}(\bu_1,\ldots,\bu_d)
	:=
	\begin{cases}
		\vol_{\HH^d}\bigl(
		S^\infty(\bu_1,\ldots,\bu_d)
		\bigr), & \ell=0, \\[1mm]
		\cH_{\HH^d}^{d-1}\bigl(
		F^\infty(\bu_1,\ldots,\bu_d)
		\bigr), & \ell=1,
	\end{cases}
\end{equation}

\begin{equation}\label{eq:C20}
	C_{2,0}=4\pi.
\end{equation}

\begin{corollary}\label{cor:LocalVolumeAsymptotics}
	Let $\ell\in\{0,1\}$ and assume $d\geq 2$ if $\ell=0$ and $d\geq 3$ if $\ell=1$. Then, as $\gamma\downarrow0$,
	\[
		\EE F^{(\ell)}_{\gamma}
		=
		C_{d,\ell}\,
		\frac{(d-2)!(d-1)^{d-2}}{d\,\kappa_{d-1}^{d-1}}\,
		\gamma
		+
		o(\gamma).
	\]
	In dimension $d=2$, the expected normalised local boundary length satisfies $\EE S_{\gamma}=2\gamma\log\frac{1}{\gamma}+O(\gamma)$.
	In particular,
	\[
		\xi_0(\lambda^{d-1},\lambda)
		=
		\frac{1}{d-1}
		-
		C_{d,0}\,
		\frac{(d-2)!(d-1)^{d-3}}{d\,\kappa_{d-1}^{d-1}}\,
		\lambda^{d-1}
		+
		o(\lambda^{d-1}).
	\]
\end{corollary}

\begin{equation}\label{eq:A}
	A(s) := \frac{\kappa_{d-1}}{d+1} s^{\frac{d+1}{2}}
		{}_2F_1\!\Big(\tfrac{d+1}{2},\tfrac{d+1}{2};\tfrac{d+3}{2};s\Big).
\end{equation}

\begin{align}
	I_d(1-t) & = C_d \, t^{-\frac{d(d-1)}{2}} \big(1 + R_I(t)\big),  \notag                                                                   \\
	A(1-t)   & = B_d \, t^{-\frac{d-1}{2}} \big(1 + R_A(t)\big), \qquad \text{with}\qquad B_d = \frac{\kappa_{d-1}}{d-1},\label{eq:Aestimate}
\end{align}

\begin{theorem}\label{thm:AsymptoticsGeneralD}
	Let $d\geq 2$ and set $\gamma=\lambda^{d-1}$. Then
	\[
		\lim_{\lambda\downarrow0}\xi_{d-1}(\lambda^{d-1},\lambda)
		=
		\frac{2^{d}}{d\,(d-1)^3}\,
		\pi^{\frac{d-2}{2}}\,
		\frac{\Gamma\!\big(\frac{(d-1)^2+1}{2}\big)}{\Gamma\!\big(\frac{(d-1)^2}{2}\big)}\,
		\left(\frac{\Gamma(\frac{d+1}{2})}{\Gamma(\frac d2)}\right)^{\,d-1}=:\xi_{d-1}^{\operatorname{PD}}\Big(\frac{1}{d-1}\Big).
	\]
\end{theorem}

\begin{proof}[Proof of Corollary~\ref{cor:LocalVolumeAsymptotics}]
	We only indicate the changes compared with the proof of
	Theorem~\ref{thm:AsymptoticsGeneralD}.
	We first consider the case \(d=2\) and \(\ell=0\). Since
	\(\cD_{\gamma,\gamma}\) is a stationary tessellation of \(\RR\) into
	intervals, its vertex and cell intensities coincide. Hence, by
	Corollary~\ref{cor:lambda_asymptotics_d2},
	\[
		\xi_0(\gamma,\gamma)
		=
		\xi_1(\gamma,\gamma)
		=
		1-\pi\gamma+o(\gamma),
		\qquad \gamma\downarrow0.
	\]
	On the other hand, Corollary~\ref{cor:Efron} and
	\(\vol_{\HH^2}(C_\gamma([0,1]))=1/\gamma\) give
	\[
		\xi_0(\gamma,\gamma)
		=
		1-\gamma\,
		\EE\vol_{\HH^2}\bigl(
		K_{\gamma,\gamma}\cap C_\gamma([0,1])
		\bigr).
	\]
	Consequently, $\EE\vol_{\HH^2}\bigl(
		K_{\gamma,\gamma}\cap C_\gamma([0,1])
		\bigr)=\pi+o(1)$,
	and therefore $\EE V_{\gamma,\gamma}([0,1])=\pi\gamma+o(\gamma)$.
	Since \(\EE V_{\gamma,\lambda}(W)\) is independent of \(W\) and
	\(\lambda\) by Theorem~\ref{thm:LocalVolume}, and since
	\(C_{2,0}=4\pi\) by \eqref{eq:C20}, this is precisely
	\[
		\EE V_{\gamma,\lambda}(W)
		=
		C_{2,0}\,
		\frac{\gamma}{4}
		+
		o(\gamma),
	\]
	as asserted. We may therefore assume \(d\ge3\) in what follows, and fix \(\ell\in\{0,1\}\).

	By Theorem~\ref{thm:LocalVolume} and
	\eqref{eqn:VolumeCylinder}, we have
	\begin{align}
		\frac{1}{\gamma}\EE F_{\gamma,\lambda}^{(\ell)}(W)
		 & =
		\frac{\gamma^{d-1}}{2d}
		\int_0^1
		(1-s)^{\frac{d-3}{2}}
		\exp\bigl(-\gamma A(s)\bigr)
		J_{d,\ell}(s)\,\dint s,
		\label{eq:LocalFunctionalAsymptotics}
	\end{align}
	where $A(s)$ is defined by \eqref{eq:A} and by  \eqref{eq:Aestimate} we have
	\[
		A(1-t)
		=
		\frac{\kappa_{d-1}}{d-1}\,
		t^{-\frac{d-1}{2}}\big(1 + R_A(t)\big),
		\qquad t\downarrow0,
	\]
	with $R_A(t) \to 0$ as $t \downarrow 0$ and $R_A$ bounded on $(0,1)$.

	It remains to determine the asymptotic behaviour of
	$J_{d,\ell}(s)$. We use the same polar substitution as in the proof of
	Theorem~\ref{thm:AsymptoticsGeneralD}, namely
	\[
		\bz_i=\sqrt{1-t y_i}\,\bu_i,
		\qquad
		\bu_i\in\SS^{d-2},
		\qquad
		y_i\in[1,1/t].
	\]
	The only additional factor, compared with $I_d(1-t)$, is
	\begin{equation}\label{eq:Auxiliary}
		\int_{
			[\sqrt{1-t y_1}\bu_1,\ldots,
					\sqrt{1-t y_d}\bu_d]}
		(1-\|\bu\|^2)^{-\frac{d-1+\ell}{2}}
		\,\dint\bu.
	\end{equation}
	This factor admits the following geometric interpretation. For
	affinely independent $\bz_1,\ldots,\bz_d\in\inter\BB^{d-1}$ put
	\[
		\bp_i
		:=
		\bigl(\bz_i,\sqrt{1-\|\bz_i\|^2}\bigr)\in\HH^d,
		\qquad i\in\{1,\ldots,d\}.
	\]
	These points lie on the Euclidean unit hemisphere
	$\Sigma:=\{(\bu,y)\in\HH^d:\|\bu\|^2+y^2=1\}$, which represents a
	totally geodesic hypersurface in the upper halfspace model. Let
	\[
		F(\bz_1,\ldots,\bz_d)
		:=
		\left\{
		\bigl(\bu,\sqrt{1-\|\bu\|^2}\bigr):
		\bu\in[\bz_1,\ldots,\bz_d]
		\right\}
	\]
	be the hyperbolic $(d-1)$-simplex in $\Sigma$ with vertices
	$\bp_1,\ldots,\bp_d$, and let
	$S(\bz_1,\ldots,\bz_d)$ be the hyperbolic $d$-simplex with facet
	$F(\bz_1,\ldots,\bz_d)$ and additional ideal vertex $\infty$, that is,
	\[
		S(\bz_1,\ldots,\bz_d)
		=
		\left\{
		(\bu,y)\in\HH^d:
		\bu\in[\bz_1,\ldots,\bz_d],\
		y\ge\sqrt{1-\|\bu\|^2}
		\right\}.
	\]
	Using the hyperbolic volume representation
	\eqref{eq:HyperbolicVolume}, we obtain
	\[
		\vol_{\HH^d}\bigl(S(\bz_1,\ldots,\bz_d)\bigr)
		=
		\frac{1}{d-1}
		\int_{[\bz_1,\ldots,\bz_d]}
		(1-\|\bu\|^2)^{-\frac{d-1}{2}}\,\dint\bu,
	\]
	while the computation of the hyperbolic area element in the proof of
	Theorem~\ref{thm:LocalVolume}, applied with $\bw=o$ and $t=1$,
	yields
	\[
		\cH_{\HH^d}^{d-1}\bigl(F(\bz_1,\ldots,\bz_d)\bigr)
		=
		\int_{[\bz_1,\ldots,\bz_d]}
		(1-\|\bu\|^2)^{-d/2}\,\dint\bu.
	\]
	Consequently, for $\bz_i=\sqrt{1-ty_i}\,\bu_i$ \eqref{eq:Auxiliary}
	equals
	\[
		\begin{cases}
			(d-1)\vol_{\HH^d}\bigl(S(\bz_1,\ldots,\bz_d)\bigr),
			 & \ell=0, \\[1mm]
			\cH_{\HH^d}^{d-1}\bigl(F(\bz_1,\ldots,\bz_d)\bigr),
			 & \ell=1.
		\end{cases}
	\]
	As $t\downarrow0$, the vertices $\bz_i=\sqrt{1-ty_i}\,\bu_i$
	converge to $\bu_i\in\SS^{d-2}$, the simplex
	$S(\bz_1,\ldots,\bz_d)$ converges to the ideal simplex
	$S^\infty(\bu_1,\ldots,\bu_d)$, and the facet
	$F(\bz_1,\ldots,\bz_d)$ converges to the ideal
	$(d-1)$-simplex $F^\infty(\bu_1,\ldots,\bu_d)$ opposite the
	vertex $\infty$. Hence, for almost every
	$(\bu_1,\ldots,\bu_d)$, this factor converges as
	$t\downarrow0$ to $(d-1)^{1-\ell}G_{d,\ell}(\bu_1,\ldots,\bu_d)$, where we recall the definition \eqref{eq:Gdl} of $G_{d,\ell}(\bu_1,\ldots,\bu_d)$. For $\ell=0$, the latter quantity is controlled by the maximal volume
	of a hyperbolic ideal $d$-simplex. For $\ell=1$ and $d\ge3$, it is
	controlled by the maximal volume of a hyperbolic ideal
	$(d-1)$-simplex. Thus the same truncation and dominated-convergence
	argument as in the proof of
	Theorem~\ref{thm:AsymptoticsGeneralD} yields
	\begin{equation}\label{eq:JdellAsymptotic}
		J_{d,\ell}(1-t)
		=
		C_{d,\ell}\,
		t^{-\frac{d(d-1)}2}
		\bigl(1+R_{d,\ell}(t)\bigr),
		\qquad t\in(0,1),
	\end{equation}
	where \(R_{d,\ell}(t)\to0\) as \(t\downarrow0\). Moreover,
	\(R_{d,\ell}\) is bounded on \((0,1)\). The above asymptotic gives
	boundedness of $R_{d,\ell}$ in a neighbourhood of zero, while away from zero this follows
	from the finiteness and monotonicity of \(J_{d,\ell}\). Namely, for every
	fixed \(t_0\in(0,1)\), we have $J_{d,\ell}(1-t)\le J_{d,\ell}(1-t_0)$, $t\in[t_0,1)$.

	To identify the constant in \eqref{eq:JdellAsymptotic}, note that the radial
	integrations contribute
	\[
		2^{-d}
		\left(
		\int_1^\infty y^{-\frac{d+1}{2}}\,\dint y
		\right)^d
		=
		\frac{1}{(d-1)^d},
	\]
	which, together with the factor $(d-1)^{1-\ell}$, gives precisely
	the normalization in the definition of $C_{d,\ell}$.

	We now make the global substitution $s=1-\gamma^{2/(d-1)}x$ with $x\in(0,\gamma^{-2/(d-1)})$.
	Using \eqref{eq:JdellAsymptotic}, all powers of $\gamma$ cancel, and
	the integrand in \eqref{eq:LocalFunctionalAsymptotics} converges
	pointwise to
	\[
		\frac{C_{d,\ell}}{2d}\,
		x^{-1-\frac{(d-1)^2}{2}}
		\exp\!\left(
		-\frac{\kappa_{d-1}}{d-1}
		x^{-\frac{d-1}{2}}
		\right).
	\]
	The boundedness of \(R_{d,\ell}\), together with the estimates for
	\(A\) used in the proof of Theorem~\ref{thm:AsymptoticsGeneralD},
	allows us to apply the same dominated-convergence argument.
	Hence
	\begin{align*}
		 & \lim_{\gamma\downarrow0}
		\frac{1}{\gamma}\,
		\EE F_{\gamma,\lambda}^{(\ell)}(W)
		=
		\frac{C_{d,\ell}}{2d}
		\int_0^\infty
		x^{-1-\frac{(d-1)^2}{2}}
		\exp\!\left(
		-\frac{\kappa_{d-1}}{d-1}
		x^{-\frac{d-1}{2}}
		\right)
		\,\dint x.
	\end{align*}
	With the substitution $u=x^{-(d-1)/2}$, the remaining integral
	equals $2\,(d-2)!\,(d-1)^{d-2}\,\kappa_{d-1}^{-(d-1)}$.
	Consequently,
	\[
		\lim_{\gamma\downarrow0}
		\frac{1}{\gamma}\,
		\EE F_{\gamma,\lambda}^{(\ell)}(W)
		=
		C_{d,\ell}\,
		\frac{(d-2)!(d-1)^{d-2}}{d\,\kappa_{d-1}^{d-1}}.
	\]
	For $\ell=0$, this yields the local-volume asymptotic, while for
	$\ell=1$ it gives the local surface-area asymptotic for $d\ge3$.
	This yields the asserted
	asymptotics for the normalised ratios
	$V_{\gamma,\lambda}(W)$ and
	$S_{\gamma,\lambda}(W)$.

	It remains to consider the local surface area in dimension $d=2$. In
	this case, a direct one-dimensional computation gives
	\begin{equation}\label{eq:J21Asymptotic}
		J_{2,1}(1-t)
		=
		4t^{-1}\log\frac{1}{t}
		+
		O(t^{-1}),
		\qquad t\downarrow0.
	\end{equation}
	Indeed, with $z_i=\frac{u_i}{\sqrt{1+u_i^2}}$ and $U=\sqrt{\frac{1-t}{t}}$,
	the integral $J_{2,1}(1-t)$ becomes
	\[
		\int_{[-U,U]^2}
		\left|
		\frac{u_1}{\sqrt{1+u_1^2}}
		-
		\frac{u_2}{\sqrt{1+u_2^2}}
		\right|
		\left|
		\operatorname{arsinh}(u_1)
		-
		\operatorname{arsinh}(u_2)
		\right|
		\,\dint u_1\dint u_2.
	\]
	The two quadrants in which $u_1$ and $u_2$ have opposite signs
	contribute
	$8U^2\log U+O(U^2)$, whereas the remaining two quadrants contribute
	only $O(U^2)$ as $U\to\infty$. Since $U^2\sim t^{-1}$ as $t\to 0$, this proves
	\eqref{eq:J21Asymptotic}.

	Substituting \eqref{eq:J21Asymptotic} into
	\eqref{eq:LocalFunctionalAsymptotics} and using
	$t=\gamma^2x$, together with
	\[
		A(1-t)=\frac{2}{3}(1-t)^{3/2}{}_2F_1\!\Big(\tfrac32,\tfrac32;\tfrac52;1-t\Big)
		=2\sqrt{\frac{1-t}{t}}-\pi+2\arcsin(\sqrt{t})
		=
		2t^{-1/2}-\pi+O(t^{1/2}),
		\quad t\downarrow0,
	\]
	gives
	\[
		\EE S_{\gamma,\lambda}(W)
		=
		2\gamma\log\frac{1}{\gamma}
		+
		O(\gamma).
	\]

	Finally, the Efron-type identity gives
	\[
		\xi_0(\lambda^{d-1},\lambda)
		=
		\frac{1}{d-1}
		-
		\lambda^{d-1}
		\EE\vol_{\HH^d}\bigl(
		K_{\lambda^{d-1},\lambda}
		\cap C_\lambda([0,1]^{d-1})
		\bigr).
	\]
	Applying the local-volume asymptotic proved above, we obtain
	\[
		\xi_0(\lambda^{d-1},\lambda)
		=
		\frac{1}{d-1}
		-
		C_{d,0}\,
		\frac{(d-2)!(d-1)^{d-3}}{d\,\kappa_{d-1}^{d-1}}\,
		\lambda^{d-1}
		+
		o(\lambda^{d-1}),
	\]
	which completes the proof.
\end{proof}

\end{document}
